% GENERATED FILE. Edit canonical lessons, then run npm run generate:books.
% canonical-source-hash: fnv1a64:be4353659ee5224e
% canonical-lessons: AR-C489-dhur, AR-C489-rub, AR-C489-hudu, AR-C489-irtaasha, AR-C489-akhafa

\chapter{Panic, terror and calm}
\label{ch:ar-panic-terror-and-calm}
\hlchaptermodality{489}

\begin{chapteropening}
\textbf{By the end of this chapter:} \emph{I can talk about panic, terror and calm, and say that someone shivered or frightened someone.}

Complete the last lesson of chapter 489: I can talk about panic, terror and calm, and say that someone shivered or frightened someone.
\end{chapteropening}

\section[\texorpdfstring{\ar{ذعر} (dhuʿr)}{ذعر (dhuʿr)}]{\ar{ذعر} (dhuʿr) — panic, alarm}

\label{lesson:AR-C489-dhur}



\begin{quote}
Before the new one: say the Arabic for to grieve, then the Arabic for to sigh.
\end{quote}

\subsection*{You'll want to know: \ar{ذعر}}
\textbf{\ar{ذعر}} — \emph{dhuʿr} — \textquotedblleft{}panic, alarm\textquotedblright{}.

\begin{quote}
\ar{أصاب} \ar{الناس} \ar{ذعر} \ar{بعد} \ar{الزلزال}.
\end{quote}

\begin{quote}
\emph{aṣāba an-nās dhuʿr baʿda az-zilzāl.}
\end{quote}

\begin{quote}
Panic seized the people after the earthquake.
\end{quote}

\begin{cousinweb}
\textbf{\ar{ذعر}} (\emph{dhuʿr}) is built on the native Arabic root \emph{dh-ʿ-r}, \textquotedblleft{}to frighten, to alarm\textquotedblright{} (Lane). Dhuʿr is a sudden fright that runs through a crowd.
\end{cousinweb}

\begin{grammarlens}[title={What you've built}]
One more word for describing what you lived and felt.
\end{grammarlens}

\subsection*{Your turn}
\begin{itemize}
\raggedright
  \item \emph{Say it:} \emph{dhuʿr}
  \item \emph{Say it:} \emph{dhuʿr}, once more
  \item \emph{Say it:} \emph{aṣāba an-nās dhuʿr baʿda az-zilzāl.} — Panic seized the people after the earthquake.
  \item \emph{Recall:} say the Arabic for to grieve, then the Arabic for to sigh, then say \emph{dhuʿr} again
\end{itemize}

\begin{tcolorbox}[breakable,colback=teal!4,colframe=teal!35!black,title={Before you move on}]
What does \ar{ذعر} mean? (\textquotedblleft{}Panic, alarm\textquotedblright{}.) Say it once more.
\end{tcolorbox}
\section[\texorpdfstring{\ar{رعب} (ruʿb)}{رعب (ruʿb)}]{\ar{رعب} (ruʿb) — terror}

\label{lesson:AR-C489-rub}



\begin{quote}
Before the new one: say the Arabic for to sigh, then the Arabic for panic.
\end{quote}

\subsection*{You'll want to know: \ar{رعب}}
\textbf{\ar{رعب}} — \emph{ruʿb} — \textquotedblleft{}terror\textquotedblright{}.

\begin{quote}
\ar{شاهد} \ar{أخي} \ar{فيلم} \ar{رعب} \ar{أمس}.
\end{quote}

\begin{quote}
\emph{shāhada akhī fīlm ruʿb ams.}
\end{quote}

\begin{quote}
My brother watched a horror film yesterday.
\end{quote}

\ar{فيلم} \ar{رعب} is a horror film.

\begin{cousinweb}
\textbf{\ar{رعب}} (\emph{ruʿb}) is built on the native Arabic root \emph{r-ʿ-b}, \textquotedblleft{}to fill with fear\textquotedblright{} (Lane). Ruʿb is fear at its strongest.
\end{cousinweb}

\begin{grammarlens}[title={What you've built}]
One more word for describing what you lived and felt.
\end{grammarlens}

\subsection*{Your turn}
\begin{itemize}
\raggedright
  \item \emph{Say it:} \emph{ruʿb}
  \item \emph{Say it:} \emph{ruʿb}, once more
  \item \emph{Say it:} \emph{shāhada akhī fīlm ruʿb ams.} — My brother watched a horror film yesterday.
  \item \emph{Recall:} say the Arabic for to sigh, then the Arabic for panic, then say \emph{ruʿb} again
\end{itemize}

\begin{tcolorbox}[breakable,colback=teal!4,colframe=teal!35!black,title={Before you move on}]
What does \ar{رعب} mean? (\textquotedblleft{}Terror\textquotedblright{}.) Say it once more.
\end{tcolorbox}
\section[\texorpdfstring{\ar{هدوء} (hudūʾ)}{هدوء (hudūʾ)}]{\ar{هدوء} (hudūʾ) — calm, quiet}

\label{lesson:AR-C489-hudu}



\begin{quote}
Before the new one: say the Arabic for panic, then the Arabic for terror.
\end{quote}

\subsection*{You'll want to know: \ar{هدوء}}
\textbf{\ar{هدوء}} — \emph{hudūʾ} — \textquotedblleft{}calm, quiet\textquotedblright{}.

\begin{quote}
\ar{أحب} \ar{هدوء} \ar{القرية} \ar{في} \ar{المساء}.
\end{quote}

\begin{quote}
\emph{uḥibbu hudūʾ al-qarya fī al-masāʾ.}
\end{quote}

\begin{quote}
I love the quiet of the village in the evening.
\end{quote}

\begin{cousinweb}
\textbf{\ar{هدوء}} (\emph{hudūʾ}) is built on the native Arabic root \emph{h-d-ʾ}, \textquotedblleft{}to become still, to be at rest\textquotedblright{} (Lane). Hudūʾ is the stillness itself: of a place, the sea or a person.
\end{cousinweb}

\begin{grammarlens}[title={What you've built}]
One more word for describing what you lived and felt.
\end{grammarlens}

\subsection*{Your turn}
\begin{itemize}
\raggedright
  \item \emph{Say it:} \emph{hudūʾ}
  \item \emph{Say it:} \emph{hudūʾ}, once more
  \item \emph{Say it:} \emph{uḥibbu hudūʾ al-qarya fī al-masāʾ.} — I love the quiet of the village in the evening.
  \item \emph{Recall:} say the Arabic for panic, then the Arabic for terror, then say \emph{hudūʾ} again
\end{itemize}

\begin{tcolorbox}[breakable,colback=teal!4,colframe=teal!35!black,title={Before you move on}]
What does \ar{هدوء} mean? (\textquotedblleft{}Calm, quiet\textquotedblright{}.) Say it once more.
\end{tcolorbox}
\section[\texorpdfstring{\ar{ارتعش} (irtaʿasha)}{ارتعش (irtaʿasha)}]{\ar{ارتعش} (irtaʿasha) — to shiver}

\label{lesson:AR-C489-irtaasha}



\begin{quote}
Before the new one: say the Arabic for terror, then the Arabic for calm.
\end{quote}

\subsection*{You'll want to know: \ar{ارتعش}}
\textbf{\ar{ارتعش}} — \emph{irtaʿasha} — \textquotedblleft{}to shiver\textquotedblright{}.

\begin{quote}
\ar{ارتعش} \ar{الولد} \ar{من} \ar{الخوف}.
\end{quote}

\begin{quote}
\emph{irtaʿasha al-walad min al-khawf.}
\end{quote}

\begin{quote}
The boy shivered with fear.
\end{quote}

Present: \ar{يرتعش}, \emph{yartaʿish}.

It is said of cold as well as fear: \ar{ارتعش} \ar{من} \ar{البرد}.

\begin{cousinweb}
\textbf{\ar{ارتعش}} (\emph{irtaʿasha}) is built on the native Arabic root \emph{r-ʿ-sh}, \textquotedblleft{}to tremble, to quiver\textquotedblright{} (Lane). Form VIII irtaʿasha is to shake all over.
\end{cousinweb}

\begin{grammarlens}[title={What you've built}]
One more word for describing what you lived and felt.
\end{grammarlens}

\subsection*{Your turn}
\begin{itemize}
\raggedright
  \item \emph{Say it:} \emph{irtaʿasha}
  \item \emph{Say it:} \emph{irtaʿasha}, once more
  \item \emph{Say it:} \emph{irtaʿasha al-walad min al-khawf.} — The boy shivered with fear.
  \item \emph{Recall:} say the Arabic for terror, then the Arabic for calm, then say \emph{irtaʿasha} again
\end{itemize}

\begin{tcolorbox}[breakable,colback=teal!4,colframe=teal!35!black,title={Before you move on}]
What does \ar{ارتعش} mean? (\textquotedblleft{}To shiver\textquotedblright{}.) Say it once more.
\end{tcolorbox}
\section[\texorpdfstring{\ar{أخاف} (akhāfa)}{أخاف (akhāfa)}]{\ar{أخاف} (akhāfa) — to frighten}

\label{lesson:AR-C489-akhafa}



\begin{quote}
Before the new one: say the Arabic for calm, then the Arabic for to shiver.
\end{quote}

\subsection*{You'll want to know: \ar{أخاف}}
\textbf{\ar{أخاف}} — \emph{akhāfa} — \textquotedblleft{}to frighten\textquotedblright{}.

\begin{quote}
\ar{أخافني} \ar{صوت} \ar{الرعد} \ar{في} \ar{الليل}.
\end{quote}

\begin{quote}
\emph{akhāfanī ṣawt ar-raʿd fī al-layl.}
\end{quote}

\begin{quote}
The sound of the thunder frightened me in the night.
\end{quote}

Present: \ar{يخيف}, \emph{yukhīf}.

Do not confuse it with \ar{أخاف} read akhāfu, I am afraid, the present of \ar{خاف}.

\begin{cousinweb}
\textbf{\ar{أخاف}} (\emph{akhāfa}) is built on the native Arabic root \emph{kh-w-f}, \textquotedblleft{}fear\textquotedblright{} (Lane). It is Form IV of khāfa, to be afraid: to make someone afraid. Mukhīf, frightening, is its participle.
\end{cousinweb}

\begin{grammarlens}[title={What you've built}]
One more word for describing what you lived and felt.
\end{grammarlens}

\subsection*{Your turn}
\begin{itemize}
\raggedright
  \item \emph{Say it:} \emph{akhāfa}
  \item \emph{Say it:} \emph{akhāfa}, once more
  \item \emph{Say it:} \emph{akhāfanī ṣawt ar-raʿd fī al-layl.} — The sound of the thunder frightened me in the night.
  \item \emph{Recall:} say the Arabic for calm, then the Arabic for to shiver, then say \emph{akhāfa} again
\end{itemize}

\begin{tcolorbox}[breakable,colback=teal!4,colframe=teal!35!black,title={Before you move on}]
What does \ar{أخاف} mean? (\textquotedblleft{}To frighten\textquotedblright{}.) Say it once more.
\end{tcolorbox}
